\section{Related Work}
\label{sec:related-work}

SPICE and its descendants~\cite{nagel1973spice} remain the reference implementation of
circuit simulation, built around modified nodal analysis~\cite{ho1975mna}. Educational tools
such as the Falstad circuit simulator and PhET's Circuit Construction Kit lower the barrier
to entry by replacing SPICE's netlist syntax with direct schematic manipulation, but all
three still present the circuit as a two-dimensional schematic. CircuitCraft uses the same
underlying numerical technique (Sections~\ref{sec:architecture}--\ref{sec:circuit-solver})
but embeds it in a three-dimensional, first-person, block-based world instead.

Minecraft has an established record as an informal STEM learning
environment~\cite{nebel2016mining,sitzmann2011meta}, with published case studies spanning
mathematics~\cite{bos2014mathematics}, general science~\cite{short2012teaching}, and
open-ended creative/collaborative use~\cite{cipollone2014creative,callaghan2016investigating}.
Within that literature, the subset closest to this paper concerns Minecraft's native redstone
subsystem, already a common entry point for teaching digital logic: Dezuanni et
al.~\cite{dezuanni2015redstone} document children reasoning about redstone circuits in
explicitly electrical language, notwithstanding that the underlying mechanic is a discrete,
boolean signal rather than a continuous one. This is the literature CircuitCraft is closest
to, and departs from most directly: those case studies use Minecraft's existing mechanics,
redstone included, as the vehicle for a lesson, whereas CircuitCraft introduces an entirely
new, continuous-time simulation layer, because redstone's discrete signal cannot represent
voltage, current, or reactive state at all.

Beyond that academic literature, several Minecraft \emph{mods} add circuit-building
mechanics of their own. Create: Circuits~\cite{chocoboyy2026createcircuits} and
ProjectRed~\cite{mrtjp2026projectred} both extend the redstone subsystem with more compact
logic-gate, arithmetic, and memory-latch blocks -- Create: Circuits labels some blocks
``analog,'' but only in the sense of Create's own multi-valued signal strength, an integer
quantity as fundamentally discrete as ordinary redstone, not a continuous electrical one;
neither mod represents voltage, current, or reactive state at all.
CircuitSim~\cite{dankov2026circuitsim} is, to the authors' knowledge, the only prior
Minecraft mod simulating genuinely analog circuits (resistors, capacitors, diodes,
transistors, op-amps), but it does so by exporting a player's circuit as a netlist and
invoking an external ngspice process installed separately on the host system, rather than
solving it natively. CircuitCraft's solver has no external dependency, is embedded directly
in the mod, and runs every game tick as an intrinsic part of the running world; to the
authors' knowledge it is the first Minecraft modification to embed a native, continuous-time
circuit solver for this purpose.

The memristor was postulated on symmetry grounds by Chua~\cite{chua1971memristor} as a
fourth fundamental circuit element; a physical device exhibiting the predicted
charge-controlled behavior was subsequently reported by Strukov et al.\ at
HP~Labs~\cite{strukov2008missing}, whose linear-drift formulation is the model implemented
here (Section~\ref{sec:circuit-solver}). Subsequent work refines this formulation for
large-scale design or device-matching purposes -- a windowed variant addressing boundary
pinning~\cite{biolek2009spice}, a threshold-based model trading physical detail for
simulation speed~\cite{kvatinsky2013team}, a physics-based compact model of filament
formation~\cite{guan2012compact}, a charge--flux-domain formulation~\cite{picos2015chargeflux},
and a simplified neuromorphic-ReRAM model~\cite{alchawa2020simple}. CircuitCraft implements
Strukov et al.'s original formulation specifically because it is simple enough to verify in
closed form (Section~\ref{sec:validation}) and traceable to the historically foundational
definition a student is most likely to encounter first, rather than optimized for matching a
fabricated device or for large-array simulation efficiency.
